\chapter{2014年北京卷（文）}

\section{单选题}

\begin{problem}
若集合 \(A = \{0, 1, 2, 4\}\)，\(B = \{1, 2, 3\}\)，则 \(A \cap B =\)（\quad）

\choices
    {\(\{0, 1, 2, 3, 4\}\)}
    {\(\{0, 4\}\)}
    {\(\{1, 2\}\)}
    {\(\{3\}\)}
\end{problem}

\begin{answer}
C
\end{answer}

\begin{solution}
集合的交集由同时属于集合 \(A\) 和集合 \(B\) 的元素组成. 集合 \(A\) 与集合 \(B\) 的公共元素为 \(1\)，\(2\)，所以 \(A\cap B=\{1,2\}\)，故选 C.
\end{solution}

\begin{problem}
下列函数中，定义域是 \(\R\) 且为增函数的是（\quad）

\choices
    {\(y = \e^{-x}\)}
    {\(y = x^{3}\)}
    {\(y = \ln x\)}
    {\(y = |x|\)}
\end{problem}

\begin{answer}
B
\end{answer}

\begin{solution}
函数 \(y=\e^{-x}\) 在 \(\R\) 上为减函数；函数 \(y=\ln x\) 的定义域为 \((0,+\infty)\)，不是 \(\R\)；函数 \(y=|x|\) 在 \((-\infty,0]\) 上为减函数，在 \([0,+\infty)\) 上为增函数，不是 \(\R\) 上的增函数. 函数 \(y=x^3\) 的定义域为 \(\R\)，且在 \(\R\) 上为增函数，所以选 B.
\end{solution}

\begin{problem}
已知向量 \(\bs{a} = (2, 4)\)，\(\bs{b} = (-1, 1)\)，则 \(2\bs{a} - \bs{b} =\)（\quad）

\choices
    {\((5,7)\)}
    {\((5,9)\)}
    {\((3,7)\)}
    {\((3,9)\)}
\end{problem}

\begin{answer}
A
\end{answer}

\begin{solution}
先计算 \(2\bs{a}=(4,8)\)，再逐坐标相减，得 \(2\bs{a}-\bs{b}=(4,8)-(-1,1)=(4+1,8-1)=(5,7)\)，故选 A.
\end{solution}

\begin{problem}
执行如图所示的程序框图，则输出的 \(S\) 值为（\quad）

\bitmapfigure[float=inline,max width=6.0cm]{img/2014/beijing_liberal/q4_fig1.png}

\choices
    {\(1\)}
    {\(3\)}
    {\(7\)}
    {\(15\)}
\end{problem}

\begin{answer}
C
\end{answer}

\begin{solution}
程序开始时 \(k=0\)，\(S=0\). 当 \(k<3\) 时，先令 \(S\leftarrow S+2^k\)，再令 \(k\leftarrow k+1\). 因此依次计算 \(k=0,1,2\) 时的三项，所得 \(S=2^0+2^1+2^2=1+2+4=7\). 此时 \(k=3\)，不再满足 \(k<3\)，程序输出 \(S=7\)，故选 C.
\end{solution}

\begin{problem}
设 \(a\)，\(b\) 是实数，则“\(a > b\)”是“\(a^2 > b^2\)”的（\quad）

\choices
    {充分且不必要条件}
    {必要且不充分条件}
    {充分必要条件}
    {既不充分也不必要条件}
\end{problem}

\begin{answer}
D
\end{answer}

\begin{solution}
条件 \(a>b\) 不能推出 \(a^2>b^2\)，例如取 \(a=1\)，\(b=-2\)，则 \(a>b\)，但 \(a^2=1<4=b^2\). 反过来，条件 \(a^2>b^2\) 也不能推出 \(a>b\)，例如取 \(a=-3\)，\(b=-2\)，则 \(a^2=9>4=b^2\)，但 \(a<b\). 因此两个条件互不充分、互不必要，故选 D.
\end{solution}

\begin{problem}
已知函数 \( f(x) = \frac{6}{x} - \log_2 x \)，在下列区间中，包含 \( f(x) \) 零点的区间是（\quad）

\choices
    {\(\left(0,1\right)\)}
    {\(\left(1,2\right)\)}
    {\(\left(2,4\right)\)}
    {\(\left(4,+\infty\right)\)}
\end{problem}

\begin{answer}
C
\end{answer}

\begin{solution}
函数 \(f(x)\) 的定义域为 \((0,+\infty)\)，且在此定义域上连续. 又因为 \(f'(x)=-\frac{6}{x^2}-\frac{1}{x\ln 2}<0\)，所以 \(f(x)\) 在 \((0,+\infty)\) 上严格递减. 计算得 \(f(2)=3-1=2>0\)，\(f(4)=\frac{3}{2}-2=-\frac{1}{2}<0\). 由零点存在定理，\(f(x)\) 在 \((2,4)\) 内有零点；由于严格递减，该零点唯一，故选 C.
\end{solution}

\begin{problem}
已知圆 \(C: (x - 3)^2 + (y - 4)^2 = 1\) 和两点 \(A(-m, 0)\)，\(B(m, 0)\)（\(m > 0\)），若圆 \(C\) 上存在点 \(P\)，使得 \(\angle APB = 90^\circ\)，则 \(m\) 的最大值为（\quad）

\choices
    {\(7\)}
    {\(6\)}
    {\(5\)}
    {\(4\)}
\end{problem}

\begin{answer}
B
\end{answer}

\begin{solution}
由圆周角定理，满足 \(\angle APB=90^\circ\) 的点 \(P\) 在以 \(AB\) 为直径的圆上. 该圆的圆心为原点，半径为 \(m\). 已知圆 \(C\) 的圆心为 \((3,4)\)，半径为 \(1\)，两圆心之间的距离为 \(\sqrt{3^2+4^2}=5\). 两圆有公共点的充要条件为 \(\lvert m-1\rvert\leqslant5\leqslant m+1\). 结合 \(m>0\)，得 \(4\leqslant m\leqslant6\)，所以 \(m\) 的最大值为 \(6\)，故选 B.
\end{solution}

\begin{problem}
加工爆米花时，爆开且不糊的粒数的百分比称为“可食用率”. 在特定条件下，可食用率 \( p \) 与加工时间 \( t \) （单位：分钟） 满足函数关系 \( p = at^2 + bt + c \) （\(a\)，\(b\)，\(c\) 是常数），下图记录了三次实验的数据. 根据上述函数模型和实验数据，可以得到最佳加工时间为（\quad）

\bitmapfigure[max width=0.32\linewidth]{img/2014/beijing_liberal/q8_fig1.png}

\choices
    {\(3.50\text{ 分钟}\)}
    {\(3.75\text{ 分钟}\)}
    {\(4.00\text{ 分钟}\)}
    {\(4.25\text{ 分钟}\)}
\end{problem}

\begin{answer}
B
\end{answer}

\begin{solution}
由图可知三次实验数据为 \(p(3)=0.7\)，\(p(4)=0.8\)，\(p(5)=0.5\)，因此
\[
\begin{cases}
9a+3b+c=0.7,\\
16a+4b+c=0.8,\\
25a+5b+c=0.5.
\end{cases}
\]
相邻两式分别相减，得 \(7a+b=0.1\)，\(9a+b=-0.3\)，两式相减得 \(2a=-0.4\)，所以 \(a=-0.2\)，进而 \(b=1.5\). 因为 \(a<0\)，抛物线开口向下，最佳加工时间为其顶点横坐标 \(t=-\dfrac{b}{2a}=-\dfrac{1.5}{2\times(-0.2)}=3.75\). 故选 B.
\end{solution}

\section{填空题}

\begin{problem}
若 \( (x+\i)\i=-1+2\i \)（\(x\in\R\)），则 \(x=\fillinblank{}\).
\end{problem}

\begin{answer}
\(2\)
\end{answer}

\begin{solution}
由 \(\i^2=-1\)，左边化简为 \((x+\i)\i=x\i+\i^2=-1+x\i\). 与右边 \(-1+2\i\) 比较虚部，得 \(x=2\).
\end{solution}

\begin{problem}
设双曲线 \(C\) 的两个焦点为 \( (- \sqrt{2}, 0) \)， \( (\sqrt{2}, 0) \)，一个顶点是 \( (1, 0) \)，则 \(C\) 的方程为\fillinblank{}.
\end{problem}

\begin{answer}
\(x^2-y^2=1\)
\end{answer}

\begin{solution}
双曲线的焦点在 \(x\) 轴上，标准方程设为 \(\dfrac{x^2}{a^2}-\dfrac{y^2}{b^2}=1\)，其中焦半距满足 \(c^2=a^2+b^2\). 由顶点为 \((1,0)\)，得 \(a=1\)；由焦点为 \((\pm\sqrt{2},0)\)，得 \(c=\sqrt{2}\). 因而 \(b^2=c^2-a^2=2-1=1\)，所以双曲线 \(C\) 的方程为 \(x^2-y^2=1\).
\end{solution}

\begin{problem}
某三棱锥的三视图如图所示，则该三棱锥的最长棱的棱长为\fillinblank{}.

\bitmapfigure[max width=0.34\linewidth]{img/2014/beijing_liberal/q11_fig1.png}
\end{problem}

\begin{answer}
\(2\sqrt{2}\)
\end{answer}

\begin{solution}
设三棱锥的顶点为 \(P\)，底面顶点为 \(A\)，\(B\)，\(C\). 由三视图可取直角坐标系，使底面位于 \(xOy\) 平面，并令
\(A=(0,0,0)\)，\(B=(2,0,0)\)，\(C=(1,1,0)\)，\(P=(0,0,2)\). 这里正视图给出高度为 \(2\) 及水平方向长度 \(2\)，俯视图给出底面三个顶点的投影为 \((0,0)\)，\((2,0)\)，\((1,1)\)，侧视图给出深度为 \(1\)，上述坐标与三视图一致.

于是六条棱长分别为
\(PA=2\)，\(PB=\sqrt{2^2+2^2}=2\sqrt{2}\)，\(PC=\sqrt{1^2+1^2+2^2}=\sqrt{6}\)，\(AB=2\)，\(AC=\sqrt{2}\)，\(BC=\sqrt{2}\). 比较它们的平方，最大值为 \(8\)，所以该三棱锥的最长棱长为 \(2\sqrt{2}\).
\end{solution}

\begin{problem}
在 \(\triangle ABC\) 中，\(a = 1\)，\(b = 2\)，\(\cos C = \frac{1}{4}\)，则 \(c =\)\fillinblank{}；\(\sin A =\)\fillinblank{}.
\end{problem}

\begin{answer}
\(2\)；\(\frac{\sqrt{15}}{8}\)
\end{answer}

\begin{solution}
由余弦定理，\(c^2=a^2+b^2-2ab\cos C=1^2+2^2-2\times1\times2\times\frac{1}{4}=4\). 因为边长为正，所以 \(c=2\).

又因为 \(C\) 是三角形内角，\(\sin C=\sqrt{1-\cos^2 C}=\sqrt{1-\frac{1}{16}}=\frac{\sqrt{15}}{4}\). 三角形 \(ABC\) 的面积既可表示为
\(\frac12ab\sin C=\frac12\times1\times2\times\frac{\sqrt{15}}4=\frac{\sqrt{15}}4\)，也可表示为 \(\frac12bc\sin A=\frac12\times2\times2\sin A=2\sin A\). 因此 \(2\sin A=\frac{\sqrt{15}}4\)，得 \(\sin A=\frac{\sqrt{15}}8\).
\end{solution}

\begin{problem}
若 \(x\)，\(y\) 满足 \(\begin{cases}
x - y - 1 \leqslant 0,\\
y \leqslant 1,\\
x + y - 1 \geqslant 0,
\end{cases}\) 则 \(z = \sqrt{3} x + y\) 的最小值为\fillinblank{}.
\end{problem}

\begin{answer}
\(1\)
\end{answer}

\begin{solution}
由约束条件得 \(x\leqslant y+1\)，\(x\geqslant1-y\)，且 \(y\leqslant1\). 要使 \(x\) 的取值区间非空，必须有 \(1-y\leqslant1+y\)，所以 \(y\geqslant0\)，从而 \(0\leqslant y\leqslant1\).

对于固定的 \(y\)，目标函数 \(z=\sqrt{3}x+y\) 随 \(x\) 增大而增大，因此取 \(x=1-y\) 时 \(z\) 最小. 此时 \(z=\sqrt{3}(1-y)+y=\sqrt{3}+(1-\sqrt{3})y\). 因为 \(1-\sqrt{3}<0\)，所以当 \(y=1\) 时取得最小值；此时 \(x=0\)，且满足原约束条件. 因而 \(z_{\min}=1\).
\end{solution}

\begin{problem}
顾客请一位工艺师把 \(A\)，\(B\) 两件玉石原料各制成一件工艺品. 工艺师带一位徒弟完成这项任务，每件原料先由徒弟完成粗加工，再由工艺师进行精加工完成制作，两件工艺品都完成后交付顾客，两件原料每道工序所需时间 （单位：工作日） 如下：

\begin{center}
\begin{tabular}{c|c|c}
\hline
原料 & 粗加工 & 精加工 \\
\hline
原料 \(A\) & \(9\) & \(15\) \\
\hline
原料 \(B\) & \(6\) & \(21\) \\
\hline
\end{tabular}
\end{center}

则最短交货期为\fillinblank{}\(\text{个工作日}\).
\end{problem}

\begin{answer}
\(42\)
\end{answer}

\begin{solution}
徒弟只能同时对一件原料进行粗加工，工艺师只能同时对一件原料进行精加工，因此只需比较徒弟加工两件原料的两种先后顺序.

若先加工原料 \(A\)，再加工原料 \(B\)，则徒弟依次在第 \(9\) 天和第 \(15\) 天完成粗加工. 工艺师从第 \(9\) 天到第 \(24\) 天完成原料 \(A\) 的精加工，随后从第 \(24\) 天到第 \(45\) 天完成原料 \(B\) 的精加工，交货期为 \(45\) 个工作日.

若先加工原料 \(B\)，再加工原料 \(A\)，则徒弟依次在第 \(6\) 天和第 \(15\) 天完成粗加工. 工艺师从第 \(6\) 天到第 \(27\) 天完成原料 \(B\) 的精加工；原料 \(A\) 虽在第 \(15\) 天完成粗加工，但要等工艺师完成原料 \(B\) 后，故从第 \(27\) 天到第 \(42\) 天完成原料 \(A\) 的精加工，交货期为 \(42\) 个工作日.

两种顺序中的最短交货期为 \(42\) 个工作日.
\end{solution}

\section{解答题}

\begin{problem}
已知 \( \{a_{n}\} \) 是等差数列，满足 \(a_{1}=3\)，\(a_{4}=12\)，数列 \( \{b_{n}\} \) 满足 \(b_{1}=4\)，\(b_{4}=20\)，且 \( \{b_{n}-a_{n}\} \) 是等比数列.

\begin{enumerate}
\item 求数列 \( \{a_{n}\} \) 和 \( \{b_{n}\} \) 的通项公式；
\item 求数列 \( \{b_{n}\} \) 的前 \(n\) 项和.
\end{enumerate}
\end{problem}

\begin{solution}
\begin{enumerate}
\item 设等差数列 \(\{a_n\}\) 的公差为 \(d\)，由题意得
\(a_4=a_1+3d\)，所以 \(12=3+3d\)，解得 \(d=3\)，从而
\(a_n=a_1+(n-1)d=3n\).

设 \(c_n=b_n-a_n\)，则 \(\{c_n\}\) 是等比数列，且
\(c_1=b_1-a_1=1\)，\(c_4=b_4-a_4=8\). 设其公比为 \(q\)，因为 \(c_1\ne0\)，所以
\(c_4=c_1q^3\)，即 \(q^3=8\)，解得 \(q=2\). 因此
\(c_n=2^{n-1}\)，于是
\(b_n=a_n+c_n=3n+2^{n-1}\).

\item 由上面的通项公式，数列 \(\{b_n\}\) 的前 \(n\) 项和为
\[
S_n=\sum_{k=1}^{n}b_k=3\sum_{k=1}^{n}k+\sum_{k=1}^{n}2^{k-1}=\frac{3n(n+1)}{2}+2^n-1
\]
所以 \(S_n=\dfrac{3n(n+1)}{2}+2^n-1\).
\end{enumerate}
\end{solution}

\begin{problem}
函数 \(f(x) = 3\sin \left(2x + \frac{\pi}{6}\right)\) 的部分图象如图所示.

\begin{enumerate}
\item 写出 \(f(x)\) 的最小正周期及图中 \(x_0\)，\(y_0\) 的值；
\item 求 \(f(x)\) 在区间 \(\left[-\frac{\pi}{2}, -\frac{\pi}{12}\right]\) 上的最大值和最小值.
\end{enumerate}

\bitmapfigure[float=inline,max width=0.62\linewidth]{img/2014/beijing_liberal/q16_fig1.png}
\end{problem}

\begin{solution}
\begin{enumerate}
\item 函数 \(f(x)=3\sin\left(2x+\dfrac{\pi}{6}\right)\) 的最小正周期为
\(T=\dfrac{2\pi}{2}=\pi\). 函数取得最大值时满足
\(2x+\dfrac{\pi}{6}=\dfrac{\pi}{2}+2k\pi\)，所以最大值为 \(3\)，从而图中的 \(y_0=3\). 第一个最大值对应的横坐标为 \(x=\dfrac{\pi}{6}\)，相邻的下一个最大值横坐标为
\(x_0=\dfrac{\pi}{6}+\pi=\dfrac{7\pi}{6}\).

\item 令 \(u=2x+\dfrac{\pi}{6}\). 当 \(x\in\left[-\dfrac{\pi}{2},-\dfrac{\pi}{12}\right]\) 时，
\[
u\in\left[-\frac{5\pi}{6},0\right]
\]
在区间 \(\left[-\dfrac{5\pi}{6},0\right]\) 上，\(\sin u\) 的最大值为 \(0\)，最小值为 \(-1\). 因此 \(f(x)\) 的最大值为 \(3\times0=0\)，最小值为 \(3\times(-1)=-3\).
\end{enumerate}
\end{solution}

\begin{problem}
如图，在三棱柱 \(ABC - A_{1}B_{1}C_{1}\) 中，侧棱垂直于底面，\(AB \perp BC\)，\(AA_{1} = AC = 2\)，\(BC = 1\)，\(E\)，\(F\) 分别为 \(A_{1}C_{1}\)，\(BC\) 的中点.

\begin{enumerate}
\item 求证：平面 \(ABE\) \(\perp\) 平面 \(B_{1}BCC_{1}\)；
\item 求证：\(C_{1}F \parallel\) 平面 \(ABE\)；
\item 求三棱锥 \(E\text{-}ABC\) 的体积.
\end{enumerate}

\bitmapfigure[float=inline,max width=0.34\linewidth]{img/2014/beijing_liberal/q17_fig1.png}
\end{problem}

\begin{solution}
\begin{enumerate}
\item 因为侧棱垂直于底面，所以 \(BB_1\perp AB\). 又因为 \(AB\perp BC\)，而直线 \(BB_1\) 与 \(BC\) 是平面 \(B_1BCC_1\) 内相交的两条直线，所以
\(AB\perp\text{平面 }B_1BCC_1\). 直线 \(AB\) 在平面 \(ABE\) 内，因此平面 \(ABE\) 垂直于平面 \(B_1BCC_1\).

\item 在底面内建立直角坐标系，并令垂直于底面的方向为 \(z\) 轴正向，取 \(B=(0,0,0)\)，\(A=(\sqrt{3},0,0)\)，\(C=(0,1,0)\). 则 \(A_1=(\sqrt{3},0,2)\)，\(C_1=(0,1,2)\). 由中点公式得 \(E=\left(\frac{\sqrt{3}}{2},\frac{1}{2},2\right)\)，\(F=\left(0,\frac{1}{2},0\right)\). 于是
\[
\overrightarrow{FC_1}=\left(0,\frac{1}{2},2\right)=\overrightarrow{BE}-\frac{1}{2}\overrightarrow{BA}.
\]
其中 \(\overrightarrow{BA}\) 与 \(\overrightarrow{BE}\) 是平面 \(ABE\) 内两条相交直线的方向向量，所以 \(\overrightarrow{FC_1}\) 是平面 \(ABE\) 内两个方向向量的线性组合. 平面 \(ABE\) 内任一点可写为
\(B+s\overrightarrow{BA}+t\overrightarrow{BE}\)，其坐标为 \(\left(\sqrt{3}s+\frac{\sqrt{3}}{2}t,\frac{1}{2}t,2t\right)\)，故平面 \(ABE\) 的方程为 \(z=4y\). 由于 \(C_1=(0,1,2)\) 不满足该方程，\(C_1\notin\) 平面 \(ABE\)，因此直线 \(C_1F\) 不在平面 \(ABE\) 内，从而 \(C_1F\parallel\) 平面 \(ABE\).

\item 在直角三角形 \(ABC\) 中，
\(AB=\sqrt{AC^2-BC^2}=\sqrt{2^2-1^2}=\sqrt{3}\)，所以
\(S_{\triangle ABC}=\dfrac{1}{2}AB\cdot BC=\dfrac{\sqrt{3}}{2}\). 点 \(E\) 到平面 \(ABC\) 的距离等于侧棱长 \(AA_1=2\)，因此
\[
V_{E\text{-}ABC}=\frac{1}{3}S_{\triangle ABC}\cdot2=\frac{\sqrt{3}}{3}
\]
所以三棱锥 \(E\text{-}ABC\) 的体积为 \(\dfrac{\sqrt{3}}{3}\).
\end{enumerate}
\end{solution}

\FloatBarrier

\begin{problem}
从某校随机抽取 100 名学生，获得了他们一周课外阅读时间（单位：小时）的数据，整理得到数据分组及频数分布表和频率分布直方图.

\begin{center}
\begin{tabular}{c|c|c}
\hline
组号 & 分组 & 频数 \\
\hline
\(1\) & \([0,2)\) & \(6\) \\
\hline
\(2\) & \([2,4)\) & \(8\) \\
\hline
\(3\) & \([4,6)\) & \(17\) \\
\hline
\(4\) & \([6,8)\) & \(22\) \\
\hline
\(5\) & \([8,10)\) & \(25\) \\
\hline
\(6\) & \([10,12)\) & \(12\) \\
\hline
\(7\) & \([12,14)\) & \(6\) \\
\hline
\(8\) & \([14,16)\) & \(2\) \\
\hline
\(9\) & \([16,18)\) & \(2\) \\
\hline
合计 & & \(100\) \\
\hline
\end{tabular}
\end{center}

\begin{enumerate}
\item 从该校随机选取一名学生，试估计这名学生该周课外阅读时间少于 12 小时的概率；
\item 求频率分布直方图中的 \(a\)，\(b\) 的值；
\item 假设同一组中的每个数据可用该组区间的中点值代替，试估计样本中的 100 名学生该周课外阅读时间的平均数在第几组.
\end{enumerate}

\bitmapfigure[float=inline,max width=0.62\linewidth]{img/2014/beijing_liberal/q20_fig1.png}
\end{problem}

\begin{solution}
\begin{enumerate}
\item 少于 12 小时的数据包括前 6 组，所以所估计的概率为
\[
\frac{6+8+17+22+25+12}{100}=\frac{90}{100}=0.9
\]
因此，这名学生该周课外阅读时间少于 12 小时的概率约为 \(0.9\).

\item 每组组距为 \(2\)，频率分布直方图的纵坐标表示频率除以组距. 因此第三组对应的高度为
\(a=\dfrac{17/100}{2}=\dfrac{17}{200}=0.085\)，第五组对应的高度为
\(b=\dfrac{25/100}{2}=\dfrac{25}{200}=0.125\).

\item 各组区间的中点依次为 \(1,3,5,7,9,11,13,15,17\)，所以样本平均数的估计值为
\[
\bar{x}=\frac{1\times6+3\times8+5\times17+7\times22+9\times25+11\times12+13\times6+15\times2+17\times2}{100}=7.68
\]
因为 \(6\leqslant7.68<8\)，所以样本平均数在第 \(4\) 组.
\end{enumerate}
\end{solution}

\begin{problem}
已知椭圆 \(C: x^{2} + 2y^{2} = 4\).

\begin{enumerate}
\item 求椭圆 \(C\) 的离心率；
\item 设 \(O\) 为原点，若点 \(A\) 在直线 \(y = 2\) 上，点 \(B\) 在椭圆 \(C\) 上，且 \(OA \perp OB\)，求线段 \(AB\) 长度的最小值.
\end{enumerate}
\end{problem}

\begin{solution}
\begin{enumerate}
\item 将椭圆方程化为标准形式
\[
\frac{x^2}{4}+\frac{y^2}{2}=1
\]
所以长半轴平方为 \(a^2=4\)，短半轴平方为 \(b^2=2\)，焦半距平方为
\(c^2=a^2-b^2=2\). 因此椭圆的离心率为
\(e=\dfrac{c}{a}=\dfrac{\sqrt{2}}{2}\).

\item 设 \(A=(u,2)\)，\(B=(x,y)\). 因为 \(B\) 在椭圆上，所以
\(x^2+2y^2=4\). 又因为 \(OA\perp OB\)，有
\(ux+2y=0\). 若 \(x=0\)，则由椭圆方程得 \(y=\pm\sqrt{2}\)，这与 \(ux+2y=0\) 矛盾，所以 \(x\ne0\)，从而 \(u=-\dfrac{2y}{x}\).

\(\triangle AOB\) 在 \(O\) 处为直角三角形，因此
\[
AB^2=OA^2+OB^2=u^2+4+x^2+y^2=4+\frac{4y^2}{x^2}+x^2+y^2
\]
令 \(s=x^2\)，则 \(0<s\leqslant4\)，且 \(y^2=\dfrac{4-s}{2}\)，所以
\[
AB^2=4+\frac{8}{s}+\frac{s}{2}
\]
由基本不等式，
\[
\frac{8}{s}+\frac{s}{2}\geqslant2\sqrt{\frac{8}{s}\cdot\frac{s}{2}}=4
\]
当且仅当 \(\dfrac{8}{s}=\dfrac{s}{2}\)，即 \(s=4\) 时取等号. 此时可取 \(B=(2,0)\)，并由垂直条件得 \(A=(0,2)\)，条件可以实现. 因而 \(AB^2\) 的最小值为 \(8\)，线段 \(AB\) 长度的最小值为 \(2\sqrt{2}\).
\end{enumerate}
\end{solution}

\begin{problem}
已知函数 \(f(x) = 2x^{3} - 3x\).

\begin{enumerate}
\item 求 \(f(x)\) 在区间 \([-2,1]\) 上的最大值；
\item 若过点 \(P(1,t)\) 存在 3 条直线与曲线 \(y = f(x)\) 相切，求 \(t\) 的取值范围；
\item 问过点 \(A(-1,2)\)，\(B(2,10)\)，\(C(0,2)\) 分别存在几条直线与曲线 \(y = f(x)\) 相切？
\end{enumerate}
\end{problem}

\begin{solution}
\begin{enumerate}
\item \(f'(x)=6x^2-3=3(2x^2-1)\)，令 \(f'(x)=0\)，得临界点
\(x=\pm\dfrac{1}{\sqrt{2}}\). 在区间 \([-2,1]\) 的端点和临界点处分别计算得
\(f(-2)=-10\)，\(f\left(-\dfrac{1}{\sqrt{2}}\right)=\sqrt{2}\)，\(f\left(\dfrac{1}{\sqrt{2}}\right)=-\sqrt{2}\)，\(f(1)=-1\).
所以 \(f(x)\) 在区间 \([-2,1]\) 上的最大值为 \(\sqrt{2}\).

\item 设切点横坐标为 \(u\). 因为 \(f'(u)=6u^2-3\)，切线方程为
\[
y=f(u)+f'(u)(x-u)=(6u^2-3)x-4u^3
\]
切线经过点 \(P(1,t)\) 的充要条件是
\[
t=-4u^3+6u^2-3=:g(u)
\]
又 \(g'(u)=-12u^2+12u=12u(1-u)\)，所以 \(g(u)\) 在 \(\left(-\infty,0\right)\) 上递减，在 \((0,1)\) 上递增，在 \(\left(1,+\infty\right)\) 上递减. 且
\(g(0)=-3\)，\(g(1)=-1\)，再结合三次函数在两端的变化趋势可知，方程 \(g(u)=t\) 有 3 个不同实根当且仅当
\(-3<t<-1\). 不同的切点横坐标对应不同的切线，因此 \(t\) 的取值范围为 \((-3,-1)\).

\item 由切线方程，过点 \((x_0,y_0)\) 的切线对应的切点横坐标 \(u\) 满足
\[
y_0=(6u^2-3)x_0-4u^3
\]
对于点 \(A(-1,2)\)，有
\[
4u^3+6u^2-1=0=\left(u+\frac{1}{2}\right)\left(4u^2+4u-2\right)
\]
其三个根 \(-\dfrac{1}{2}\)，\(\dfrac{-1+\sqrt{3}}{2}\)，\(\dfrac{-1-\sqrt{3}}{2}\) 互不相同，所以过点 \(A\) 有 3 条切线.

对于点 \(B(2,10)\)，有
\[
4u^3-12u^2+16=4(u-2)^2(u+1)=0
\]
其不同实根为 \(u=2,-1\)，所以过点 \(B\) 有 2 条切线.

对于点 \(C(0,2)\)，有
\(2=-4u^3\)，即 \(u^3=-\dfrac{1}{2}\)，只有 1 个实根，所以过点 \(C\) 有 1 条切线.

不同的 \(u\) 不会给出同一条切线，否则同一条直线会与三次曲线在两个不同点处均相切，交点重数将超过三次. 因此结论为：过点 \(A\)、\(B\)、\(C\) 分别有 \(3\)、\(2\)、\(1\) 条切线.
\end{enumerate}
\end{solution}
